%% tables/tbl_roadmap.tex  (개선 로드맵 — 개선 논리를 한곳에 가시화, 상태=TBD)
%% 각 lever: 타깃 병목 → 개선 논리(메커니즘) → 기대 효과 → 상태.

\begin{table*}[t]
\centering
\caption{\textbf{Host-side reclamation 개선 로드맵.}
  각 lever 가 어떤 host 병목을 어떤 논리로 줄이는지와 현재 상태.
  정량 결과는 \S\ref{subsec:decisive} 결정적 실험 후 채운다(상태 \textbf{TBD}).
  C 는 물리적으로 패배하여 폐기.}
\label{tbl:roadmap}
\footnotesize
\begin{tabular}{@{}l >{\raggedright\arraybackslash}p{0.17\textwidth} >{\raggedright\arraybackslash}p{0.40\textwidth} >{\raggedright\arraybackslash}p{0.20\textwidth} l@{}}
\toprule
\textbf{Lever} & \textbf{타깃 host 병목} & \textbf{개선 논리 (메커니즘)} & \textbf{기대 효과} & \textbf{상태} \\
\midrule
A1 CPU 드래프팅 & draft 제안(suffix tree/소형 모델) & 드래프트 생성을 유휴 CPU 코어에 병렬화해 GPU verify 뒤로 숨김 → GPU 는 verify 만 수행 & GPU util 회복, suffix 이득 확대 & \textbf{TBD} \\
A2 KV tiering & HBM 한계로 배치/컨텍스트 제약 & 유휴 DRAM 2TB(+CXL)로 cold KV 계층화 → in-flight 시퀀스$\uparrow$ → GPU gap 을 다른 요청이 충전 & 유효 배치$\uparrow$, util 회복 & \textbf{TBD} \\
B1 detokenize 오프로드 & 고-tps detok 직렬화 & 토큰 단위 detok 을 별도 CPU 코어로 분리·병렬화 → step 경로에서 제거 & feeding 지연 단축 & \textbf{TBD} \\
B2 constrained-decode & 문법/JSON 마스크·FSM 계산 & step 마다 발생하는 mask 생성을 CPU 로 오프로드 (code/tool 워크로드 집중) & code 워크로드 stall 제거 & \textbf{TBD} \\
B3 스케줄러/샘플링 파이프라인 & step 간 스케줄·샘플 직렬 비용 & host 파이프라인을 다중 코어로 분할해 GPU feeding 과 겹침 & GPU 공급률$\uparrow$ & \textbf{TBD} \\
\midrule
\textit{C} GPU 매트멀 오프로드 & --- & CPU 가 GPU 연산을 대신 (H2D/D2H 동기 비용) & 음(陰) & \textbf{폐기} \\
\bottomrule
\end{tabular}
\end{table*}
